% =====================================================================
\chapter{Cardinalidad, palomar y diagonal de Cantor}
\label{tema:6}
% =====================================================================

Las funciones inyectivas, sobreyectivas y biyectivas permiten comparar el tamaño de dos conjuntos sin contar sus elementos. En conjuntos finitos esta comparación da el principio del palomar; en conjuntos infinitos conduce a la distinción entre conjuntos numerables y no numerables y al argumento diagonal de Cantor. El tema sigue la sección~12.3 y el capítulo~14 de Hammack~\cite{hammack}, el capítulo~8 de Velleman~\cite{velleman}, las secciones~4.5, 7.1 y~14.8 de~\cite{mcs} y las transparencias «Cardinality» de Stanford CS103~\cite{cs103card}.

\section{Comparar tamaños con funciones}

\begin{definicion}[Equipotencia]
Dos conjuntos $A$ y $B$ tienen el \emph{mismo cardinal}, y se escribe $\abs{A}=\abs{B}$, si existe una biyección $A\to B$. Se escribe $\abs{A}\le\abs{B}$ si existe una función inyectiva $A\to B$.
\end{definicion}

Para conjuntos finitos, estas definiciones coinciden con contar elementos (\cref{prop:t3-finitos}). La relación «tener el mismo cardinal» es reflexiva ($\id_A$ es biyectiva), simétrica (la inversa de una biyección es biyectiva, \cref{prop:t4-inv-comp}) y transitiva (la composición de biyecciones es biyectiva, \cref{teo:t4-comp}).

\begin{ejemplo}[Un conjunto infinito y una parte propia]\label{ej:t6-pares}
Sea $2\N=\{0,2,4,\dots\}$. La función $f\colon\N\to2\N$, $f(n)=2n$, es biyectiva: es inyectiva porque $2n=2m$ implica $n=m$, y sobreyectiva porque cada par $2k$ es $f(k)$. Por tanto $\abs{\N}=\abs{2\N}$ aunque $2\N\subsetneq\N$. Para conjuntos finitos esto es imposible: una parte propia tiene estrictamente menos elementos.
\end{ejemplo}

\section{El principio del palomar}

\begin{teorema}[Principio del palomar]\label{teo:t6-palomar}
Sean $A$ y $B$ finitos y $f\colon A\to B$.
\begin{enumerate}
  \item Si $\abs{A}>\abs{B}$, $f$ no es inyectiva: existen $a\neq a'$ con $f(a)=f(a')$.
  \item (Forma general.) Si $\abs{A}>k\abs{B}$ con $k\in\N$, alguna fibra tiene al menos $k+1$ elementos. En particular, alguna fibra tiene al menos $\lceil\abs{A}/\abs{B}\rceil$ elementos.
\end{enumerate}
\end{teorema}

\begin{proof}
(ii) Si todas las fibras tuvieran a lo sumo $k$ elementos, se tendría $\abs{A}=\sum_{b\in B}\abs{f^{-1}(\{b\})}\le k\abs{B}$, en contra de la hipótesis. El apartado~(i) es el caso $k=1$. Para la última afirmación, sea $m$ el tamaño de la mayor fibra; entonces $\abs{A}\le m\abs{B}$, luego $m\ge\abs{A}/\abs{B}$ y, por ser entero, $m\ge\lceil\abs{A}/\abs{B}\rceil$.
\end{proof}

El nombre procede de la imagen de $\abs{A}$ palomas que se refugian en $\abs{B}$ palomares: si hay más palomas que palomares, en algún palomar hay al menos dos. La dificultad en las aplicaciones está en elegir bien las «palomas» y los «palomares».

\begin{ejemplo}[Restos]\label{ej:t6-restos}
Entre cinco enteros cualesquiera hay dos cuya diferencia es múltiplo de~4. Palomas: los cinco enteros; palomares: los cuatro restos posibles al dividir entre~4. Dos enteros con el mismo resto tienen diferencia múltiplo de~4.
\end{ejemplo}

\begin{proposicion}[Límite de la compresión sin pérdida]\label{prop:t6-compresion}
Para $n\ge1$ no existe ninguna función inyectiva de $\{0,1\}^{n}$ en el conjunto $\{0,1\}^{<n}=\bigcup_{k=0}^{n-1}\{0,1\}^{k}$ de cadenas binarias de longitud menor que $n$.
\end{proposicion}

\begin{proof}
$\abs{\{0,1\}^{n}}=2^{n}$, mientras que $\abs{\{0,1\}^{<n}}=\sum_{k=0}^{n-1}2^{k}=2^{n}-1$. Por el \cref{teo:t6-palomar}(i), ninguna función entre ellos es inyectiva.
\end{proof}

Un compresor sin pérdida es una función inyectiva (debe poder descomprimirse). La \cref{prop:t6-compresion} dice que ningún compresor acorta \emph{todas} las cadenas de longitud $n$: si acorta algunas, alarga o deja igual otras. El módulo~6 cuantifica esta idea con la entropía.

\section{Conjuntos numerables}

\begin{definicion}[Numerable]
Un conjunto $A$ es \emph{infinito numerable} si $\abs{A}=\abs{\N}$, es decir, si sus elementos pueden listarse como $a_0,a_1,a_2,\dots$ sin repeticiones y sin omitir ninguno. Es \emph{numerable} si es finito o infinito numerable, y \emph{no numerable} en otro caso.
\end{definicion}

\begin{proposicion}\label{prop:t6-numerables}
Los conjuntos $\Z$, $\N\times\N$ y $\Sigma^{*}$ (con $\Sigma$ finito y no vacío) son infinitos numerables.
\end{proposicion}

\begin{proof}
\emph{El conjunto $\Z$.} La función del \cref{eje:3-NZ} es una biyección $\N\to\Z$, que da la lista $0$, $-1$, $1$, $-2$, $2$, \dots

\emph{El conjunto $\N\times\N$.} Se recorren los pares por diagonales $m+n=0,1,2,\dots$, y dentro de cada diagonal por orden creciente de $n$. La diagonal $m+n=d$ tiene $d+1$ pares, y las anteriores contienen $1+2+\dots+d=d(d+1)/2$; por tanto, el par $(m,n)$ ocupa la posición
\[
\pi(m,n)=\frac{(m+n)(m+n+1)}{2}+n.
\]
Cada natural $N$ se alcanza exactamente una vez: hay un único $d$ con $d(d+1)/2\le N<(d+1)(d+2)/2$, y entonces $n=N-d(d+1)/2\in\{0,\dots,d\}$ y $m=d-n$. Luego $\pi$ es una biyección (la \emph{función de emparejamiento de Cantor}).

\emph{El conjunto $\Sigma^{*}$.} Se listan las cadenas por longitud creciente y, dentro de cada longitud, en orden lexicográfico (\emph{orden de longitud y lexicográfico}). Cada longitud aporta un número finito $\abs{\Sigma}^{T}$ de cadenas, de modo que toda cadena aparece en una posición finita.
\end{proof}

\begin{figure}[h]
\centering
\begin{tikzpicture}[scale=0.8,font=\small,>={Stealth[length=1.6mm]}]
  \foreach \m in {0,...,3} \foreach \n in {0,...,3} {
    \pgfmathtruncatemacro{\d}{\m+\n}
    \pgfmathtruncatemacro{\p}{\d*(\d+1)/2+\n}
    \node[circle,draw,inner sep=1.5pt,minimum size=6mm] (p\m\n) at (\n*1.3,-\m*1.0) {\p};}
  \foreach \m in {0,...,3} \node[anchor=east] at (-0.55,-\m*1.0) {$m=\m$};
  \foreach \n in {0,...,3} \node at (\n*1.3,0.8) {$n=\n$};
  \draw[->,gray] (p10) -- (p01);
  \draw[->,gray] (p20) -- (p11); \draw[->,gray] (p11) -- (p02);
  \draw[->,gray] (p30) -- (p21); \draw[->,gray] (p21) -- (p12); \draw[->,gray] (p12) -- (p03);
\end{tikzpicture}
\caption{Recorrido de $\N\times\N$ por diagonales. Cada círculo muestra la posición $\pi(m,n)$ del par $(m,n)$.}
\label{fig:t6-cantor}
\end{figure}

\begin{proposicion}\label{prop:t6-criterios}
Sea $A$ un conjunto no vacío. Son equivalentes: (i) $A$ es numerable; (ii) existe una función inyectiva $A\to\N$; (iii) existe una función sobreyectiva $\N\to A$. Además, todo subconjunto de un conjunto numerable es numerable, y la unión de una familia numerable de conjuntos numerables es numerable.
\end{proposicion}

Se admite este resultado sin demostración completa (véase~\cite[§14.2]{hammack} y~\cite[§8.2]{velleman}); algunas de sus partes se piden en los \cref{eje:6-subconjN,eje:6-producto}. Como consecuencia, $\Q$ es numerable (\cref{eje:6-Q}).

\section{El argumento diagonal de Cantor}

\begin{teorema}[Cantor]\label{teo:t6-cantor}
Para todo conjunto $A$ no existe ninguna función sobreyectiva $A\to\Pa(A)$. En consecuencia, $\abs{A}\neq\abs{\Pa(A)}$.
\end{teorema}

\begin{proof}
Sea $f\colon A\to\Pa(A)$ una función cualquiera. Se considera el conjunto «diagonal»
\[
D=\{x\in A: x\notin f(x)\}\subseteq A.
\]
Supóngase que $D=f(y)$ para algún $y\in A$. Si $y\in D$, por la definición de $D$, $y\notin f(y)=D$; si $y\notin D$, entonces $y\notin f(y)$ y, por la definición de $D$, $y\in D$. Ambos casos son contradictorios, luego $D$ no es imagen de ningún elemento y $f$ no es sobreyectiva.
\end{proof}

\begin{corolario}\label{cor:t6-nonum}
$\Pa(\N)$ y el conjunto $\{0,1\}^{\N}$ de sucesiones infinitas de ceros y unos no son numerables.
\end{corolario}

\begin{proof}
Por el \cref{teo:t6-cantor}, no hay biyección $\N\to\Pa(\N)$, y $\Pa(\N)$ es infinito; luego no es numerable. La función $X\mapsto\chi_X$ es una biyección $\Pa(\N)\to\{0,1\}^{\N}$ (\cref{tema:2}), así que $\{0,1\}^{\N}$ tampoco es numerable.
\end{proof}

El nombre «diagonal» se entiende con sucesiones: dada una lista $s_0,s_1,s_2,\dots$ de sucesiones binarias, la sucesión $d$ con $d_k=1-s_k(k)$ difiere de $s_k$ en la posición $k$, luego no está en la lista. El mismo argumento, con desarrollos decimales y evitando las colas de nueves, prueba que $\R$ no es numerable~\cite[§14.2]{hammack}.

\begin{observacion}[Conexión con el curso]
El conjunto $\V^{*}$ de todos los textos tokenizados es infinito numerable (\cref{prop:t6-numerables}); por eso un modelo de lenguaje es una distribución \emph{discreta} y la condición de normalización es una serie (módulo~1, §1.3; módulo~5). El principio del palomar explica por qué un identificador de token de 15~bits no basta para un vocabulario de $5\cdot10^{4}$ tokens (módulo~12) y por qué toda función \emph{hash} sobre un conjunto mayor que su rango produce colisiones (módulo~16). Por último, los números en coma flotante forman un conjunto \emph{finito} (módulo~11): una red implementada en un ordenador es una función entre conjuntos finitos, aunque se modele con funciones reales.
\end{observacion}

\begin{observacion}[Errores frecuentes]
(1) Pensar que una parte propia de un conjunto infinito es siempre «más pequeña» (\cref{ej:t6-pares}). (2) Creer que todos los conjuntos infinitos tienen el mismo cardinal: $\abs{\N}<\abs{\Pa(\N)}$. (3) Aplicar el palomar sin identificar con precisión la función de las palomas a los palomares. (4) En el argumento diagonal con decimales, olvidar que $0{,}0999\ldots=0{,}1$: hay que elegir dígitos que eviten esa ambigüedad.
\end{observacion}

\begin{lecturas}
\item Hammack~\cite{hammack}, §3.9 «The Division and Pigeonhole Principles», §12.3 «The Pigeonhole Principle Revisited» y cap.~14 «Cardinality of Sets», §14.1–14.3.
\item Velleman~\cite{velleman}, §8.1 «Equinumerous Sets» y §8.2 «Countable and Uncountable Sets».
\item Lehman, Leighton y Meyer~\cite{mcs}, §4.5 «Finite Cardinality», §7.1 «Infinite Cardinality» y §14.8 «The Pigeonhole Principle».
\item Stanford CS103~\cite{cs103card}, transparencias «Cardinality».
\end{lecturas}

\section*{Ejercicios del tema 6}
\addcontentsline{toc}{section}{Ejercicios del tema 6}

\begin{ejercicio}{1}{eje:6-vf}
Decida si son verdaderas o falsas, justificando con un resultado del tema: (a) si $A\subsetneq B$, entonces $\abs{A}<\abs{B}$; (b) $\abs{\N}=\abs{2\N}$; (c) $\Q$ es numerable; (d) $\Pa(\N)$ es numerable; (e) existe una biyección entre $\Z$ y $\N\times\N$.
\end{ejercicio}

\begin{ejercicio}{1}{eje:6-textos}
¿Hay más textos posibles (elementos de $\V^{*}$, con $\V$ finito) que números naturales? ¿Más que números reales? ¿Y más que números en coma flotante de 32~bits? Justifique cada respuesta.
\end{ejercicio}

\begin{ejercicio}{2}{eje:6-meses}
(a) En un grupo de 13 personas, ¿hay necesariamente dos nacidas el mismo mes? (b) ¿Cuántas personas hacen falta para garantizar que tres hayan nacido el mismo mes? (c) Se extraen calcetines de un cajón con calcetines de 4 colores; ¿cuántos hay que extraer para garantizar un par del mismo color?
\end{ejercicio}

\begin{ejercicio}{2}{eje:6-bits}
(a) Un vocabulario tiene $K=5\cdot10^{4}$ tokens. Demuestre que los identificadores no caben en 15~bits y calcule el número mínimo de bits necesario. (b) Se asignan $10^{6}$ palabras a $2^{16}$ cubetas de una tabla \emph{hash}. Demuestre que alguna cubeta recibe al menos 16 palabras.
\end{ejercicio}

\begin{ejercicio}{3}{eje:6-restos}
(a) Demuestre que entre $n+1$ enteros cualesquiera hay dos cuya diferencia es múltiplo de $n$. (b) Demuestre que si se eligen $n+1$ números del conjunto $\{1,2,\dots,2n\}$, dos de ellos son consecutivos.
\end{ejercicio}

\begin{ejercicio}{3}{eje:6-puntos}
Demuestre que entre cinco puntos cualesquiera del plano con coordenadas enteras hay dos cuyo punto medio tiene también coordenadas enteras. ¿Es cierto con cuatro puntos?
\end{ejercicio}

\begin{ejercicio}{3}{eje:6-compresion}
(a) Demuestre la \cref{prop:t6-compresion}. (b) Sea $C\colon\{0,1\}^{*}\to\{0,1\}^{*}$ una función inyectiva. Demuestre que para cada $n\ge1$ existe una cadena $w$ de longitud $n$ con $\abs{C(w)}\ge n$.
\end{ejercicio}

\begin{ejercicio}{2}{eje:6-pi}
Con la función de emparejamiento $\pi(m,n)=\tfrac{(m+n)(m+n+1)}{2}+n$: (a) calcule $\pi(2,3)$ y $\pi(0,4)$; (b) halle el par $(m,n)$ con $\pi(m,n)=17$; (c) compruebe en la \cref{fig:t6-cantor} los valores de la diagonal $m+n=3$.
\end{ejercicio}

\begin{ejercicio}{2}{eje:6-shortlex}
Sea $\Sigma=\{\cad{a},\cad{b}\}$ con $\cad{a}<\cad{b}$. (a) Escriba las 14 primeras cadenas de $\Sigma^{*}$ en el orden de longitud y lexicográfico, numeradas desde~0. (b) ¿En qué posición aparece $\cad{bab}$? (c) Demuestre que, si se lee $\cad{a}$ como 0 y $\cad{b}$ como 1, la cadena de longitud $L$ cuyo valor binario es $v$ ocupa la posición $2^{L}-1+v$.
\end{ejercicio}

\begin{ejercicio}{2}{eje:6-diagonal}
Las cinco primeras sucesiones binarias de una supuesta lista de $\{0,1\}^{\N}$ empiezan así: $s_0=10110\ldots$, $s_1=00101\ldots$, $s_2=11100\ldots$, $s_3=01011\ldots$, $s_4=11111\ldots$. Escriba los cinco primeros términos de la sucesión diagonal $d$, con $d_k=1-s_k(k)$, y compruebe que $d$ difiere de cada $s_k$, $k\le4$.
\end{ejercicio}

\begin{ejercicio}{3}{eje:6-subconjN}
Demuestre que todo subconjunto infinito $S\subseteq\N$ es infinito numerable: construya una biyección $\N\to S$ que liste los elementos de $S$ en orden creciente.
\fuente{adaptado de Cambridge, Discrete Mathematics~\cite{camdm}, hoja~8, ejercicio~13.2.1.}
\end{ejercicio}

\begin{ejercicio}{3}{eje:6-producto}
(a) Demuestre que $g\colon\N\times\N\to\N$, $g(m,n)=2^{m}(2n+1)-1$, es biyectiva. (b) Deduzca que si $A$ y $B$ son infinitos numerables, $A\times B$ también lo es.
\fuente{adaptado de Cambridge, Discrete Mathematics~\cite{camdm}, hoja~8, ejercicio~13.2.3.}
\end{ejercicio}

\begin{ejercicio}{3}{eje:6-Q}
Demuestre que $\Q$ es numerable. Puede usar que todo racional se escribe de forma única como $p/q$ con $p\in\Z$, $q\in\N$, $q\ge1$ y $\operatorname{mcd}(\abs{p},q)=1$, y la \cref{prop:t6-criterios}.
\fuente{adaptado de Cambridge, Discrete Mathematics~\cite{camdm}, hoja~8, ejercicio~13.1.2.}
\end{ejercicio}

\begin{ejercicio}{3}{eje:6-PA}
Sea $A$ un conjunto cualquiera. Demuestre que no existe ninguna función inyectiva $\Pa(A)\to A$.
\fuente{adaptado de Cambridge, Discrete Mathematics~\cite{camdm}, hoja~8, ejercicio~13.2.4.}
\end{ejercicio}

\begin{ejercicio}{3}{eje:6-intervalos}
(a) Demuestre que $f\colon[0,1]\to[0,2]$, $f(x)=2x$, es biyectiva. (b) Construya una biyección $[0,1]\to(0,1)$ y demuestre que lo es. (Indicación: envíe $0\mapsto\tfrac12$ y $\tfrac1n\mapsto\tfrac1{n+2}$ para $n\ge1$, y deje fijos los demás puntos.)
\fuente{adaptado de Stanford CS103~\cite{cs103ps3}, Problem Set~3, problema~7 «Infinity Minus Two».}
\end{ejercicio}
